\subsection{形式幂级数的定义。阶}

设 I 为一个集合。我们回顾（III，第 454 与 456 页），A 上的幺半群 \(\mathbf{N}^{(I)}\) 的总体代数称为关于未定元 \(\mathbf{X}_i\) ( \(i \in I\) )（或在未定元 \(\mathbf{X}_i\) 中）的、系数在 A 中的形式幂级数代数。它记作 \(\mathbf{A}[[\mathbf{X}_i]]_{i \in I}\) 或 \(\mathbf{A}[[(\mathbf{X}_i)_{i \in I}]]\) 或也记作 \(\mathbf{A}[[\mathbf{X}]]\) ，这时以 \(\mathbf{X}\) 表示族 \((\mathbf{X}_i)_{i \in I}\) ：在本节中我们将主要使用记号 \(\mathbf{A}[[I]]\) 。有时方便的是，在 \(\mathbf{A}[[I]]\) 中把 I 的元素 \(i\) 的典范像用异于 \(\mathbf{X}_i\) 的符号来表示，例如 \(\mathbf{Y}_i\) , \(\mathbf{Z}_i\) , \(\mathrm{T}_i\) , \ldots；这种情形下所用的约定与多项式情形的约定类似（IV，第 1 页）。代数 \(\mathbf{A}[[I]]\) 这时记作 \(\mathbf{A}[[\mathbf{Y}_i]]_{i \in I}\) ，或 \(\mathbf{A}[[\mathbf{Y}]]\) 等等。

当 I 是含 p 个元素的有限集合时，我们也说 A{[}{[}I{]}{]} 是 p 个未定元的形式幂级数代数。对固定的 p，这些代数都是同构的。1, 2, \ldots{} 个未定元的形式幂级数代数也将记作 A{[}{[}X{]}{]}, A{[}{[}U, V{]}{]}, \ldots，而不具体指出指标集 I。

形式幂级数 \(u\) 通常写作 \(u = \sum_{\nu \in \mathbb{N}^{(I)}} \alpha_{\nu} X^{\nu}\) （参见 IV，第 1 页）。

这些 \(\alpha_{\nu}\) 是 \(u\) 的系数；它们之中可以有无限多个 \(\neq 0\) 。这些 \(\alpha_{\nu}X^{\nu}\) 称为 \(u\) 的项； \(u\) 是多项式的充分必要条件是 \(u\) 只有有限个 \(\neq 0\) 的项。这些项 \(\alpha_{\nu}X^{\nu}\) （满足 \(|\nu|=p\) ）称为总次数为 p 的项。形式幂级数 \(u_{p}=\sum_{|\nu|=p}\alpha_{\nu}X^{\nu}\) 称为 u 的 p 次齐次分量（当 I 有限时它是一个多项式）； \(u_{0}\) 等同于 A 的一个元素，也称为 u 的常数项。我们说 u 是 p 次齐次的，如果 \(u = u_{p}\) 。若 u， \(v \in A[[I]]\) 且 w = uv，我们有

\[
w _ {p} = \sum_ {q + r = p} u _ {q} v _ {r}\tag{1}
\]

对每个整数 \(p \geqslant 0\) 成立。

我们回顾（III，第 456 页），阶 \(\omega(u)\) （形式幂级数 \(u \neq 0\) 的）是最小整数 \(p\) （满足 \(u_p \neq 0\) ）。我们约定在 \(\mathbf{Z}\) 中添加一个记作 \(\infty\) 的元素，并按下列约定把 \(\mathbf{Z}\) 上的序关系与加法延拓到 \(\mathbf{Z} \cup \{\infty\}\) ：

\[
n <   \infty , \quad \infty + \infty = \infty , \quad \infty + n = n + \infty = \infty
\]

对每个 \(n \in \mathbf{Z}\) ；我们还设 \(\omega(0) = \infty\) 。在这些约定下，我们有关系式

\[
\begin{array}{c} \omega (u + v) \geqslant \inf (\omega (u), \omega (v)), \\ \omega (u + v) = \inf (\omega (u), \omega (v)) \quad \text { if } \quad \omega (u) \neq \omega (v), \\ \omega (u v) \geqslant \omega (u) + \omega (v), \end{array}
\]

（对任意形式幂级数 \(u\) 与 \(v\) （在 A{[}{[}I{]}{]} 中）成立。）

我们回顾（III，第 457 页），对任意子集 \(J\) （ \(I\) 的），我们把 \(A[[I]]\) 与 \(A[[I - J]][[J]]\) 等同，这使我们能够定义阶 \(\omega_{j}(u)\) （形式幂级数相对于 \(X_{j}\) ( \(j \in J\) ) 的）、 \(u\) 相对于 \(X_{j}\) ( \(j \in J\) ) 的齐次分量等等。

设 \(\varphi\) 是 A 到环 B 的一个同态。我们把 \(\varphi\) 延拓为同态 \(\overline{\varphi}\) （ A{[}{[}I{]}{]} 到 B{[}{[}I{]}{]} 的），其方式是让每个形式幂级数 \(u = \sum_{\nu} \alpha_{\nu} X^{\nu}\) 对应

到形式幂级数 \(\sum_{\nu} \varphi(\alpha_{\nu}) X^{\nu}\) ；我们说后者是通过

将 \(\varphi\) 应用到形式幂级数 \(u\) 的系数上而得到的。我们有时用 \({}^{\varphi}u\) 代替 \(\bar{\varphi}(u)\) 。

特别地，如果 A 是 B 的子环且 \(\varphi\) 是 A 到 B 的典范单射，则同态 \(\bar{\varphi}\) （ A{[}{[}I{]}{]} 到 B{[}{[}I{]}{]} 的）是单射；我们通常把 A{[}{[}I{]}{]} 与 B{[}{[}I{]}{]} 的一个子环等同起来（借助 \(\bar{\varphi}\) ）。

